\scanpage{190}
[N.B. $\operatorname{Pic}(\widetilde C-X)\simeq\operatorname{Pic}(\unsure{\bar C|\widetilde C})/\mathbb Z.\unsure{\eta|y}$, où où $\eta$ est la classe de $X$ considéré comme diviseur dans $\widetilde C$ ; \deleted{\unsure{$\xi$|et}} Comme $\operatorname{Pic}(\widetilde C)\xleftarrow{\sim}\operatorname{Pic}(X)$ par $\widetilde C\mathbin{-}X$, on retrouve ainsi l'isomorphisme $\operatorname{Pic}(C-a)\simeq\operatorname{Pic}(X)/\mathbb Z.\eta$].

D'autre part, on trouve aussi
\[
\operatorname{Pic}(\widehat S-\hat a)\simeq\operatorname{Pic}(\widetilde{\widehat S}-X)\simeq\operatorname{Pic}(\widetilde{\widehat S})/\mathbb Z\hat\eta
\]
où $\widetilde{\widehat S}$ désigne l'éclaté de $\widehat S$. Or on a un isom formel évident
\[
\widetilde{\widehat S}\simeq\widetilde S\times_S\widehat S\ \underset{\smash{\tikz[baseline=0pt]{\draw (0,.08)--(0,0)--(1.05,0)--(1.05,.08);}}}{\simeq\widehat{\widetilde S}}
\quad\text{(complété formel de $\widetilde S$ le long de $X$)}
\]
d'où, par le th. d'existence
\[
\operatorname{Pic}(\widetilde{\widehat S})\simeq\operatorname{Pic}(\widehat{\widetilde S})\ ;
\]
\deleted{comme par ailleurs} Ainsi, le \deleted{$\operatorname{Pic}(\widetilde{\widehat S}\unsure{-\hat a}$} morphisme
\[
\operatorname{Pic}(\deleted{\unsure{C_a}}S-a)\mathbin{-}\operatorname{Pic}(\widehat{\deleted{\unsure{C_a}}S}-a)
\]
est induit, en passant au quotient par $\mathbb Z\eta$ resp. $\mathbb Z\hat\eta$, par le morphisme
\[
\operatorname{Pic}(E)\longrightarrow\operatorname{Pic}(\widehat E)
\]
% archive sheet number: 189
